% LENS B — Robotics / Automation roles
% (Senior Mechatronics Engineer, Robotics Engineer, Automation Engineer)
\documentclass[10pt,a4paper]{article}
\usepackage[margin=1.6cm]{geometry}
\usepackage{enumitem}
\usepackage{textcomp}
\setlist[itemize]{leftmargin=*, itemsep=1pt, topsep=2pt, parsep=0pt}
\usepackage{titlesec}
\titleformat{\section}{\large\bfseries}{}{0pt}{}[\vspace{-4pt}\rule{\textwidth}{0.4pt}]
\titlespacing*{\section}{0pt}{8pt}{4pt}
\pagestyle{empty}

\newcommand{\Headline}{Senior Mechatronics \& Robotics Engineer}

\newcommand{\Summary}{Senior Mechatronics Engineer with 6+ years' experience designing, building, and deploying automated electro-mechanical systems --- from closed-loop test rigs integrating motors, sensors, DAQ, and embedded controllers, to a field-deployed robotic coating-application system that turned a manual industrial process into a repeatable automated workflow. I work across the full stack of physical automation: mechanism and system design, device integration and I/O, control and scripting (Python, C++, ROS/ROS2), calibration and error handling in real operating environments, and validation through structured test campaigns to production-grade robustness. Currently completing a PhD in Biomedical Engineering / Collaborative Robotics (physics-based simulation of human movement --- MATLAB, OpenSim, MuJoCo).}

\newcommand{\Competencies}{Robotic system design, build \& deployment \; $\cdot$ \; Device \& instrument integration (motors, actuators, ToF/current/voltage sensing, DAQ, embedded controllers) \; $\cdot$ \; Python \& C++ for device control \& workflow automation \; $\cdot$ \; ROS/ROS2 \; $\cdot$ \; Closed-loop control, calibration \& error handling \; $\cdot$ \; Translating manual processes into automated workflows \; $\cdot$ \; Physics-based simulation (MATLAB/Simulink, MuJoCo, OpenSim, Gazebo) \; $\cdot$ \; Motion planning, kinematics \& PID control \; $\cdot$ \; Test-campaign data logging \& analysis \; $\cdot$ \; Prototype-to-production}

\newcommand{\AssetCoolBullets}{
  \item Led the design, build, and deployment of the Dynamic Coating Application Module (DCAM) --- a robotic system automating the application of protective coatings onto power-grid conductors, carried on base robotic platforms --- taking it from proof-of-concept prototype to a robust, repeatable, field-deployed system.
  \item Translated a manual coating process into an automated workflow in close collaboration with chemists and fluidics specialists: capturing process requirements (delivery rate, uniformity, coverage), designing the actuation and fluid-delivery system to meet them, and building in the sensing and controls needed for repeatable, in-spec application.
  \item Integrated the module's devices end-to-end --- actuators, motor drivers, ToF/current/voltage sensing, and embedded control --- working with the electronics and software teams on I/O, device behaviour, calibration routines, and error handling for safe, reliable operation in a harsh outdoor environment; ran structured lab test campaigns across build iterations (3D-printed concept to CNC-machined production-intent), logging test and sensor data to verify performance against specification.
}

\newcommand{\EducationExtra}{}

\newcommand{\Skills}{%
\textbf{Robotics \& Controls:} Forward/inverse kinematics \& dynamics, motion planning, trajectory generation, PID control, closed-loop systems, calibration \& error handling\\
\textbf{Software \& Scripting:} Python (NumPy), C++, MATLAB, ROS/ROS2, Git, real-time systems, data analysis, applied machine learning\\
\textbf{Simulation \& Modelling:} MATLAB/Simulink, Robotics System Toolbox, MuJoCo, OpenSim, Gazebo, Unity, FEA, Motor-CAD (electromagnetic/thermal)\\
\textbf{Device Integration \& Electronics:} Sensors and DAQ, ToF/current/voltage sensing, BLDC/FOC motor control, motor drivers, embedded controllers, CAN protocol work, signal conditioning, safety circuits, PCB collaboration\\
\textbf{Mechanical \& Prototyping:} SolidWorks, Onshape, Fusion 360 (all Advanced); DFMA, tolerance analysis, precision mechanisms, rapid prototyping (3D printing, CNC, WaterJet)\\
\textbf{Data \& Test:} Test-campaign design and execution, instrumented data logging, performance characterisation and validation against specification\\
\textbf{Collaboration:} Cross-disciplinary work with scientists (chemistry, fluidics), electronics, software, and customer stakeholders; junior engineer mentoring; global/matrix teams}

\newcommand{\Interests}{Open-source robotics, embedded systems projects, hardware prototyping, applied AI for physical systems, mentoring engineers.}

\begin{document}
% ============================================================
% SPINE — shared, canonical content. Facts never change here.
% Rulings: decommissioned line (never "live"/"live-line"/"energised");
% canonical titles fixed; only lens files vary emphasis.
% ============================================================

\begin{center}
  {\LARGE\bfseries Uljan Sinani}\\[2pt]
  {\large \Headline{} $|$ United Kingdom}\\[3pt]
  uljan.sinani@gmail.com \; $\cdot$ \; 073-7977-8240 \; $\cdot$ \; linkedin.com/in/uljansinani \; $\cdot$ \; Full UK driving licence\\[2pt]
  {\small Portfolio: \textbf{sinani.ai} — physics-based simulation, world-model latent audit, open engineering notebook \; $\cdot$ \; github.com/USinani}
\end{center}

\section*{Professional Summary}
\Summary

\section*{Core Competencies}
\Competencies

\section*{Professional Experience}

\noindent\textbf{AssetCool Ltd., Leeds, United Kingdom} \hfill May 2025 -- June 2026\\
\textit{Senior Mechatronics Engineer (Mechanical Design Lead)}
\begin{itemize}
  \AssetCoolBullets
  \item Delivered international field commissioning and validation with Hydro-Qu\'ebec, Canada (June 2025) --- full uniform 360\textdegree{} coating across a 300\,m conductor span on a decommissioned line --- including installation health \& safety planning for site work on utility grid infrastructure.
  \item Coordinated a multidiscipline design effort across mechanical, electronics, software, chemistry, and fluidics teams; held the mechanical design accountable to functional requirements and client specification, managing design change between prototype and production-intent builds.
\end{itemize}

\noindent\textbf{Federal-Mogul Powertrain / Tenneco, Essex} \hfill 2023 -- 2025\\
\textit{Mechatronics Design Engineer (Lead)}
\begin{itemize}
  \item Led a five-person team investigating and redesigning the power-electronics water cooling plate on the Belt Integrated Starter-Generator (AMG programme): hydraulic/thermal design, tolerance analysis, and FEA, delivered as a complete 3D model and 2D drawing package to BS~8888 with GD\&T; coordinated the checking, approval, and validation handover (shaker and integration testing) against customer specification.
  \item Owned the mechanical design of the generator housing through a client-driven change programme (new connectors and mounting points for range expansion) --- implementing design change control and reporting progress directly to the end customer alongside internal stakeholders.
  \item Built an electromagnetic/thermal simulation of the switched reluctance motor in Motor-CAD to characterise flux and heat dissipation; used the temperature-gradient insight to define optimal thermistor placement and advise the power-electronics team on power modulation and conductor/harness routing to minimise EMI.
  \item Designed, prototyped, and validated a precision thermistor clamp for repeatable sensor placement within the stator bore, carrying the design from lab verification through to final customer approval at the Munich production site.
  \item Delivered technical drawings to relevant ISO and BS standards for the housing, cooling plate, and rotor; reviewed and approved design work from junior engineers and mentored the team on analysis methods, CAD standards, and documentation practice.
\end{itemize}

\noindent\textbf{Borg Automotive UK Ltd., West Midlands} \hfill 2019 -- 2020\\
\textit{Mechatronics Engineer (R\&D)}
\begin{itemize}
  \item Designed bespoke electro-mechanical test fixtures and automated validation systems for electric power-steering assemblies --- closed-loop test rigs integrating motors, sensors, DAQ, and embedded controllers; conducted stress, thermal, and dynamic FEA to verify sub-assembly robustness under operational loads.
  \item Standardised CAD modelling conventions and 2D drawing standards across the design office; conducted design reviews to improve documentation quality, consistency, and checking efficiency.
\end{itemize}

\noindent\textbf{Car Parts Industries (CPI), Belgium} \hfill 2018 -- 2019\\
\textit{Mechatronics Engineer -- R\&D}
\begin{itemize}
  \item R\&D and New Product Introduction for electric steering columns (VW, FIAT, Opel/Vauxhall); on the Prometheus project, developed an embedded module interfacing with the CAN communication between steering units and the vehicle ECU --- hands-on work with automotive connectors, harnessing, and electrical interconnects.
\end{itemize}

\noindent\textbf{BORG Automotive A/S, Denmark} \hfill 2016 -- 2017\\
\textit{Junior Mechatronics Engineer}
\begin{itemize}
  \item New Product Introduction for electric and mechanical steering racks; authored test-rig specifications, developed jigs and diagnostic test equipment, and fed structured issue reports back to the remanufacturing team.
\end{itemize}

\section*{Doctoral Research}
\noindent\textbf{University of Reading --- PhD, Biomedical Engineering / Collaborative Robotics (distance, part-time)} \hfill Nov 2021 -- Present\\
\textit{Physics-based biomechanical simulation of human upper-limb motion}
\begin{itemize}
  \item Developing two- and three-link rigid-body dynamic models of the human arm in MATLAB with anthropometric parameterisation; integrating Hill-type muscle actuators (OpenSim) and building 3D visualisation pipelines in MuJoCo and Unity --- applying structured modelling, simulation, and performance-evaluation methods to complex mechanical system behaviour.
\end{itemize}

\section*{Education}
\noindent\textbf{Tallinn University of Technology --- MSc, Mechatronics Engineering} \hfill 2014 -- 2016\\
Exchange, University of Southern Denmark --- Mechatronics for Nanotechnology\EducationExtra\\
\textbf{Polytechnic University of Tirana --- BSc, Mechatronics Engineering} \hfill 2012 -- 2014

\section*{Technical Skills}
\Skills

\section*{Courses \& Interests}
Udacity Robotics Software Engineer Nanodegree \; $|$ \; Industrial Automation \; $|$ \; Leadership in Engineering\\
\Interests

\end{document}
